\thesisabstract{%
  This thesis is a series of publications that introduce novel methods for neural rendering using limited information, focusing on Neural Radiance Fields (NeRFs) and 3D Gaussian Splatting (3DGS).
  It explores how these models construct 3D representations from 2D images and
  demonstrates ways to condition these representations for generating
  high-quality renderings of 3D objects.

  We present that both NeRFs and 3DGS representations are not
  \textit{controllable} in an interpretable way---once trained, one can no
  longer change the content of the scene such that output reflects the changes
  intended in the input.
  We propose four complementary techniques that use simple, interpretable inputs
  derived from sparse training data and extend these methods to perform
  effectively in few-shot learning scenarios.

  We begin by examining the area of Neural Radiance Fields addressing
  limitations in existing approaches and presenting contributions to achieve
  controllable radiance fields.
  By incorporating partial and sparse data during training, our NeRF
  extension---\conerf---leverages the smoothness of neural networks to produce
  controllable, high-quality images across general scene types.
  We argue that sparse, manual annotations obtainable in the matter of a minute
  are sufficient to provide an interpretable control over the NeRF's output
  thanks to the spectral bias of neural networks.
  We support the argument with extensive evaluation results, showcasing that
  \conerf is among the first steps to fully controllable radiance fields.

  To tackle the reliance on extensive, high-quality data annotations from
  multi-view videos, we introduce a new method, \blendfields, for training
  neural radiance fields in few-shot, multi-view settings.
  This approach learns internal deformation templates, which blend smoothly
  during inference, significantly improving image quality compared to existing
  baselines and enabling effective rendering from limited input images.
  We leverage our volume representation in the form of tetrahedral cages.
  Using traditional techniques in computer graphics, we compute physically-based
  deformation quantities that provide smooth and interpretable interpolation
  between known expressions.
  We show in our experiments that \blendfields achieves higher quality of
  renderings from just a few shots which contrasts to previous works that rely
  on hundreds of frames in a single capture.

  Moreover, we develop a novel model, \lumigauss, for controlling radiance
  fields through environmental lighting for ``in-the-wild'' captures under
  varying daylight conditions.
  By incorporating precomputed radiance transfer well-known in computer
  graphics, this model enables physically plausible scene relighting and
  provides users with intuitive control over lighting in reconstructed scenes.
  We achieve this goal despite the fact that each training image is captured
  under different lighting conditions and does not contain any information about
  light presented in the scene.

  Many applications of neural rendering require high performance, thus in the
  last work of this thesis we focus on adaptable computation efficiency with our
  \clog.
  It applies a coarse-to-fine learning strategy for 3D Gaussian Splatting, which
  upscales a 2D grid that stores Gaussian latent representations.
  This strategy achieves competitive results while allowing deployment on
  various computational devices with minimal quality loss.
  As demonstrated in our experiments, \clog becomes increasingly efficient
  during inference with fewer Gaussians.

  This research advances the state of the art in controllable neural radiance
  fields and expands their application to few-shot learning scenarios.
  These innovations enhance the possibilities for neural rendering and other
  objects in the scene from limited information and open new directions for
  future research in the field.
}

\thesiskeywords{Neural Rendering, Neural Radiance Fields, Few-Shot Learning, Human Rendering, Partial Information, Gaussian Splatting}

